\section{Turing degree and a quartic normal form}
\label{red:degree-normal-form}

The finite-test reduction determines the Turing degree of rational
solvability. A separate arithmetic-circuit conversion gives the same
degree for two restricted forms of the input.

\begin{corollary}\label{thm:degree-normal-form}
With the usual effective coding, let \(\mathrm{H10}(\Q)\) be the set of
integral polynomials having a rational zero, with the number of variables
part of the input. Then \(\mathrm{H10}(\Q)\) has Turing degree \(0'\),
the degree of the halting problem. The same is true for each of the
following restricted input problems:
\begin{enumerate}[label=(\roman*),leftmargin=*]
\item rational solvability for integral polynomials of total degree at
most four;
\item given a nonempty finite list
\(q_1,\ldots,q_s\in\Z[X_1,\ldots,X_N]\), each of total degree at most
two, whether \(\sum_{i=1}^s q_i^2\) has a zero in \(\Q^N\).
\end{enumerate}
No bound is imposed on the number of variables, the number of summands
in (ii), or the coefficient heights.
\end{corollary}

\begin{proof}
Let \(\mathrm{H10}(\Z)\) denote the analogous set of integral polynomials
having an integer zero. The parallel search in the proof of Theorem~\ref{thm:main}
in Section~\ref{sec:reduction} can be run with
an \(\mathrm{H10}(\Q)\) oracle in place of the hypothetical rational
decision procedure. The finite tests of \S\ref{red:finite-tests} are
uniformly computable from \(f\) and their indices once the fixed
arithmetic data have been chosen, and each test uses finitely many oracle
queries. The search therefore halts on every input, proving
\[
                    \mathrm{H10}(\Z)\le_{\mathrm T}\mathrm{H10}(\Q),
\]
where \(\le_{\mathrm T}\) denotes Turing reducibility. The fixed choices
used to define \(\mathcal T\) are independent of \(f\), so this
construction gives a single oracle machine.

The Davis--Putnam--Robinson--Matiyasevich theorem, in the form
\cite[Theorem~8.1]{Davis1973}, applies to a positive-integer coding of
the halting set. Its positive-integer witnesses can be replaced by
\(1+u_1^2+u_2^2+u_3^2+u_4^2\), using the four-square theorem as in
\cite[p.~234]{Davis1973}. Specializing the parameter in the resulting
fixed polynomial gives a computable many-one reduction of the halting
set to \(\mathrm{H10}(\Z)\). Integer solvability is itself computably
enumerable by enumeration of integer tuples, so it has Turing degree
\(0'\). Rational solvability is also computably enumerable: enumerate
rational tuples and evaluate the input polynomial exactly. Every
computably enumerable set is Turing reducible to the halting set.
Together with the displayed oracle reduction, this proves that
\(\mathrm{H10}(\Q)\) has degree \(0'\).

For the degree restriction, we apply over \(\Q\) the algebraic
normalization used in \cite[Theorem~7.5]{Davis1973}. Given
\(f\in\Z[X_1,\ldots,X_n]\), effectively construct an arithmetic circuit
evaluating \(f\), with input nodes, integral constant nodes, and addition
and multiplication gates. Give every node \(v\) a new variable \(Y_v\).
For an input node \(X_j\) and a constant node \(c\), impose respectively
\[
                         Y_v-X_j=0,\qquad Y_v-c=0.
\]
For a sum or product node \(v\) with incoming nodes \(u,w\), impose
respectively
\[
                    Y_v-Y_u-Y_w=0,\qquad Y_v-Y_uY_w=0.
\]
Finally impose \(Y_{\mathrm{out}}=0\) at the output node. This gives a
nonempty finite list \(q_1,\ldots,q_s\) of integral polynomials of total
degree at most two in \(N=n+m\) common variables
\(\mathbf Z=(\mathbf X,\mathbf Y)\), where \(m\) is the number of circuit
nodes. It also handles constant inputs. Induction through the circuit shows
that each rational input tuple has a unique extension satisfying the
gate equations, with output value \(f(\mathbf X)\). Since squares of
rationals are nonnegative, it follows that
\[
\begin{gathered}
 (\exists\mathbf x\in\Q^n)\ f(\mathbf x)=0
 \\
 \Longleftrightarrow\quad
 (\exists\mathbf z\in\Q^N)\ \bigwedge_{i=1}^s \bigl(q_i(\mathbf z)=0\bigr)\\
 \Longleftrightarrow\quad
 (\exists\mathbf z\in\Q^N)\ \sum_{i=1}^s q_i(\mathbf z)^2=0.
\end{gathered}
\]
Let \(S_2\) denote the presented list problem in (ii), and let
\(H_{\le4}\) denote the problem in (i). Sending \(f\) to the list
\((q_1,\ldots,q_s)\), and then expanding its squared sum, gives computable
many-one reductions
\[
                    \mathrm{H10}(\Q)\le_{\mathrm m}S_2
                                      \le_{\mathrm m}H_{\le4}.
\]
The expanded polynomial has integral coefficients and total degree at
most four. Both target sets are computably enumerable by rational-tuple
enumeration, so the reductions and the first part of the proof give
Turing degree \(0'\) for both.

The many-one reduction to the restricted problems starts from
\(\mathrm{H10}(\Q)\); the reduction from \(\mathrm{H10}(\Z)\) remains a
Turing reduction. This argument does not establish many-one completeness
of \(\mathrm{H10}(\Q)\). The list in (ii) is supplied as part
of the input, so no recognition of sum-of-squares presentations is
needed. The construction supplies no bound on the number of variables
or coefficient heights, and it gives no single existential definition
of \(\Z\) in \(\Q\).
\end{proof}
